% !TeX root = ../../paper.tex

\subsection{Passive Haptics and Proxy Objects}
\label{sec:rw-passive-haptics}

The core observation behind passive haptics predates consumer VR by two decades. Hoffman showed in 1998 that letting participants pick up a real plate while seeing its virtual twin made them rate other virtual objects as more solid and the environment as more realistic~\cite{Hoffman.1998}. Insko's dissertation work made the case more dramatically: participants who felt a low wooden ledge under their feet while standing at a virtual pit showed stronger presence responses than those who only saw the drop~\cite{Insko.2001}. Both studies used proxies that were far simpler than the objects they stood in for, which is what makes the technique attractive---the haptic channel apparently does not demand the fidelity the visual channel does.

How much discrepancy players tolerate is the question Simeone et al.\ addressed with Substitutional Reality~\cite{substitutional-reality}. They systematically mismatched physical objects and their virtual representations---an umbrella standing in for a lightsaber being the canonical example---and derived guidelines for which properties (size, shape, temperature, weight) can drift and which break the illusion. Their participants kept accepting substitutions well past the point of literal resemblance, provided the mismatch did not contradict what the hand felt at first contact. Azmandian et al.\ pushed in the other direction with haptic retargeting, warping the virtual world so that a single physical cube could serve as the proxy for several virtual objects in different locations~\cite{Azmandian.2016}.

Our project sits deliberately at the low-tech end of this spectrum. We use no retargeting and no instrumented props; each virtual object has one fixed physical counterpart chosen for its texture rather than its geometry. What distinguishes our setup from the studies above is the emotional register of the materials. Prior proxy work has mostly paired neutral objects---cubes, plates, furniture. We pair proxies with content designed to disgust, which raises a question the literature leaves open: whether confirming the player's worst tactile expectations amplifies the emotional response, or whether the mundane identity of the proxy (a gummy worm is, after all, candy) undercuts it.
